% Options for packages loaded elsewhere
\PassOptionsToPackage{unicode}{hyperref}
\PassOptionsToPackage{hyphens}{url}
\PassOptionsToPackage{dvipsnames,svgnames,x11names}{xcolor}
%
\documentclass[
  11pt,
  a4paper,
]{article}

\usepackage{amsmath,amssymb}
\usepackage{iftex}
\ifPDFTeX
  \usepackage[T1]{fontenc}
  \usepackage[utf8]{inputenc}
  \usepackage{textcomp} % provide euro and other symbols
\else % if luatex or xetex
  \usepackage{unicode-math}
  \defaultfontfeatures{Scale=MatchLowercase}
  \defaultfontfeatures[\rmfamily]{Ligatures=TeX,Scale=1}
\fi
\usepackage{lmodern}
\ifPDFTeX\else  
    % xetex/luatex font selection
    \setmainfont[]{Inter}
\fi
% Use upquote if available, for straight quotes in verbatim environments
\IfFileExists{upquote.sty}{\usepackage{upquote}}{}
\IfFileExists{microtype.sty}{% use microtype if available
  \usepackage[]{microtype}
  \UseMicrotypeSet[protrusion]{basicmath} % disable protrusion for tt fonts
}{}
\makeatletter
\@ifundefined{KOMAClassName}{% if non-KOMA class
  \IfFileExists{parskip.sty}{%
    \usepackage{parskip}
  }{% else
    \setlength{\parindent}{0pt}
    \setlength{\parskip}{6pt plus 2pt minus 1pt}}
}{% if KOMA class
  \KOMAoptions{parskip=half}}
\makeatother
\usepackage{xcolor}
\setlength{\emergencystretch}{3em} % prevent overfull lines
\setcounter{secnumdepth}{5}
% Make \paragraph and \subparagraph free-standing
\makeatletter
\ifx\paragraph\undefined\else
  \let\oldparagraph\paragraph
  \renewcommand{\paragraph}{
    \@ifstar
      \xxxParagraphStar
      \xxxParagraphNoStar
  }
  \newcommand{\xxxParagraphStar}[1]{\oldparagraph*{#1}\mbox{}}
  \newcommand{\xxxParagraphNoStar}[1]{\oldparagraph{#1}\mbox{}}
\fi
\ifx\subparagraph\undefined\else
  \let\oldsubparagraph\subparagraph
  \renewcommand{\subparagraph}{
    \@ifstar
      \xxxSubParagraphStar
      \xxxSubParagraphNoStar
  }
  \newcommand{\xxxSubParagraphStar}[1]{\oldsubparagraph*{#1}\mbox{}}
  \newcommand{\xxxSubParagraphNoStar}[1]{\oldsubparagraph{#1}\mbox{}}
\fi
\makeatother


\providecommand{\tightlist}{%
  \setlength{\itemsep}{0pt}\setlength{\parskip}{0pt}}\usepackage{longtable,booktabs,array}
\usepackage{calc} % for calculating minipage widths
% Correct order of tables after \paragraph or \subparagraph
\usepackage{etoolbox}
\makeatletter
\patchcmd\longtable{\par}{\if@noskipsec\mbox{}\fi\par}{}{}
\makeatother
% Allow footnotes in longtable head/foot
\IfFileExists{footnotehyper.sty}{\usepackage{footnotehyper}}{\usepackage{footnote}}
\makesavenoteenv{longtable}
\usepackage{graphicx}
\makeatletter
\newsavebox\pandoc@box
\newcommand*\pandocbounded[1]{% scales image to fit in text height/width
  \sbox\pandoc@box{#1}%
  \Gscale@div\@tempa{\textheight}{\dimexpr\ht\pandoc@box+\dp\pandoc@box\relax}%
  \Gscale@div\@tempb{\linewidth}{\wd\pandoc@box}%
  \ifdim\@tempb\p@<\@tempa\p@\let\@tempa\@tempb\fi% select the smaller of both
  \ifdim\@tempa\p@<\p@\scalebox{\@tempa}{\usebox\pandoc@box}%
  \else\usebox{\pandoc@box}%
  \fi%
}
% Set default figure placement to htbp
\def\fps@figure{htbp}
\makeatother

\makeatletter
\@ifpackageloaded{caption}{}{\usepackage{caption}}
\AtBeginDocument{%
\ifdefined\contentsname
  \renewcommand*\contentsname{Table of contents}
\else
  \newcommand\contentsname{Table of contents}
\fi
\ifdefined\listfigurename
  \renewcommand*\listfigurename{List of Figures}
\else
  \newcommand\listfigurename{List of Figures}
\fi
\ifdefined\listtablename
  \renewcommand*\listtablename{List of Tables}
\else
  \newcommand\listtablename{List of Tables}
\fi
\ifdefined\figurename
  \renewcommand*\figurename{Figure}
\else
  \newcommand\figurename{Figure}
\fi
\ifdefined\tablename
  \renewcommand*\tablename{Table}
\else
  \newcommand\tablename{Table}
\fi
}
\@ifpackageloaded{float}{}{\usepackage{float}}
\floatstyle{ruled}
\@ifundefined{c@chapter}{\newfloat{codelisting}{h}{lop}}{\newfloat{codelisting}{h}{lop}[chapter]}
\floatname{codelisting}{Listing}
\newcommand*\listoflistings{\listof{codelisting}{List of Listings}}
\makeatother
\makeatletter
\makeatother
\makeatletter
\@ifpackageloaded{caption}{}{\usepackage{caption}}
\@ifpackageloaded{subcaption}{}{\usepackage{subcaption}}
\makeatother

\usepackage{bookmark}

\IfFileExists{xurl.sty}{\usepackage{xurl}}{} % add URL line breaks if available
\urlstyle{same} % disable monospaced font for URLs
\hypersetup{
  pdftitle={Discussion},
  pdfauthor={Nambona YANGUERE},
  colorlinks=true,
  linkcolor={\#33FFA2},
  filecolor={Maroon},
  citecolor={\#FF33FF},
  urlcolor={Blue},
  pdfcreator={LaTeX via pandoc}}


\title{Discussion}
\usepackage{etoolbox}
\makeatletter
\providecommand{\subtitle}[1]{% add subtitle to \maketitle
  \apptocmd{\@title}{\par {\large #1 \par}}{}{}
}
\makeatother
\subtitle{Institutional Resilience and the Medieval ESG Parallel}
\author{Nambona YANGUERE}
\date{2026-04-03}

\begin{document}
\maketitle

\renewcommand*\contentsname{Table of contents}
{
\hypersetup{linkcolor=}
\setcounter{tocdepth}{3}
\tableofcontents
}

\section{1. Orders as Institutional
Protocols}\label{orders-as-institutional-protocols}

The central finding of this study is that \textbf{the religious order
functions as a deterministic institutional protocol} that shapes
economic output, territorial strategy, and survivability.

This is not a metaphor. The Rule of St.~Benedict, adopted by both
Benedictine and Cistercian communities, functioned as a
\textbf{governance framework} comparable to modern corporate charters:

\begin{itemize}
\tightlist
\item
  \textbf{Resource allocation rules} --- hours of labor, agricultural
  quotas
\item
  \textbf{Succession protocols} --- election of abbots, term limitations
\item
  \textbf{Compliance mechanisms} --- visitation systems, chapter
  oversight
\item
  \textbf{Expansion strategy} --- filiation charters for new foundations
\end{itemize}

The data confirms that communities adhering strictly to codified rules
show measurably higher governance scores and longer institutional
lifespans. This is consistent with modern institutional economics:
\textbf{codified governance reduces agency risk and improves
resilience.}

\section{2. Land Transformation as Environmental
Capital}\label{land-transformation-as-environmental-capital}

\subsection{2.1 The Polder Case}\label{the-polder-case}

The development of polders in Flanders and Northern France represents
one of the most significant land transformation projects in European
history. Cistercian communities were primary actors in this
transformation, converting tidal marshlands into productive agricultural
land.

From an ESG perspective, this activity presents a dual reading:

\begin{itemize}
\tightlist
\item
  \textbf{Positive:} Creation of durable agricultural infrastructure
  that increased regional food security and economic output for
  centuries beyond the communities' own lifespans
\item
  \textbf{Contextual:} This transformation altered natural ecosystems. A
  modern ESG assessment would require a biodiversity impact evaluation
  that historical records do not support
\end{itemize}

This study scores land transformation on the basis of \textbf{documented
economic output and institutional intent}, acknowledging that
retrospective environmental scoring carries inherent methodological
limitations.

\subsection{2.2 Viticulture and Agricultural
Diversification}\label{viticulture-and-agricultural-diversification}

Monastic communities were responsible for systematic viticulture across
Burgundy, the Rhine Valley, and the Iberian Peninsula. The data shows
that communities with \textbf{diversified agricultural output} (multiple
crop types + livestock + viticulture) score 12 points higher on the
environmental dimension than monoculture-dependent sites.

This finding parallels modern ESG literature on agricultural portfolio
diversification as a resilience strategy.

\section{3. Geopolitical Positioning and
Resilience}\label{geopolitical-positioning-and-resilience}

\subsection{3.1 Trade Route Proximity}\label{trade-route-proximity}

Communities located within \textbf{30 km of a documented trade route}
show:

\begin{itemize}
\tightlist
\item
  15\% higher revenue indices
\item
  20\% longer institutional lifespans
\item
  Lower variance in economic output across centuries
\end{itemize}

This suggests that \textbf{connectivity to commercial infrastructure}
was a primary determinant of institutional sustainability, independent
of order affiliation.

\subsection{3.2 Frontier vs.~Interior
Positioning}\label{frontier-vs.-interior-positioning}

Cistercian frontier communities demonstrate a higher-risk, higher-reward
profile:

\begin{itemize}
\tightlist
\item
  Higher peak economic output during stable periods
\item
  Sharper declines during geopolitical disruption
\item
  Greater dependence on agricultural self-sufficiency
\end{itemize}

Benedictine interior communities show lower volatility, benefiting from
diversified revenue sources (education, hospitality, urban proximity).

This mirrors the modern financial distinction between \textbf{growth
assets} (high-return, high-volatility) and \textbf{defensive assets}
(stable-return, low-volatility).

\section{4. The Dissolution Problem}\label{the-dissolution-problem}

\subsection{4.1 Political Intervention as Systemic
Risk}\label{political-intervention-as-systemic-risk}

The dominant cause of dissolution (42\%) is \textbf{political
intervention}: royal seizures, reformation-era suppressions, and
territorial redistribution.

This finding positions political risk as the \textbf{primary threat to
institutional continuity}, regardless of internal governance quality.
Communities with high governance scores were not immune to external
seizure, indicating that institutional resilience has structural limits
when facing sovereign-level actors.

\subsection{4.2 Economic Collapse as Liquidity
Risk}\label{economic-collapse-as-liquidity-risk}

The second cause (28\%) is economic failure. Communities dependent on a
single revenue source (typically agricultural output) were most
vulnerable.

This parallels modern portfolio theory: \textbf{concentration risk
reduces institutional resilience.} Communities with diversified revenue
(agriculture + education + trade + patronage) survived economic shocks
at twice the rate of single-source communities.

\subsection{4.3 Governance Failure}\label{governance-failure}

Internal governance failure accounts for 12\% of dissolutions. These
cases correlate with:

\begin{itemize}
\tightlist
\item
  Extended leadership vacancies
\item
  Documented violations of the founding rule
\item
  Conflict between community and diocesan authority
\end{itemize}

This confirms that \textbf{governance protocol adherence} is a
measurable predictor of institutional health, consistent with modern ESG
governance scoring.

\section{5. Relevance to Modern ESG
Frameworks}\label{relevance-to-modern-esg-frameworks}

\subsection{5.1 What Medieval Data Teaches Modern
Scoring}\label{what-medieval-data-teaches-modern-scoring}

This study demonstrates that the core principles of modern ESG analysis
--- environmental stewardship, social contribution, and governance
quality --- have measurable historical precedents.

The medieval monastic system offers a \textbf{longitudinal dataset}
spanning 800+ years that modern ESG frameworks lack. Corporate ESG
scoring typically covers 5--20 years of data. Monastic records allow us
to test whether ESG-aligned behaviors produce durable institutional
outcomes over centuries.

The answer from this dataset is \textbf{conditionally yes}:

\begin{itemize}
\tightlist
\item
  High governance scores correlate with longevity
\item
  Diversified environmental strategies correlate with economic
  resilience
\item
  Social contribution (education, hospitality) correlates with network
  density
\end{itemize}

The condition is that \textbf{no ESG score protects against
sovereign-level political intervention}, a finding relevant to modern
impact investors operating in jurisdictions with political instability.

\subsection{5.2 Heritage as Long-Term
Impact}\label{heritage-as-long-term-impact}

A dimension absent from standard ESG scoring is \textbf{cultural and
architectural heritage}. Monastic communities produced:

\begin{itemize}
\tightlist
\item
  Architectural innovations (Gothic, Romanesque)
\item
  Agricultural techniques (polders, viticulture, crop rotation)
\item
  Educational infrastructure (scriptoria, libraries, schools)
\item
  Landscape transformation visible to this day
\end{itemize}

These contributions persist centuries after institutional dissolution.
Modern ESG frameworks, focused on quarterly or annual reporting, do not
capture this category of \textbf{ultra-long-term impact}. This study
suggests that a ``Heritage'' dimension could enrich modern impact
scoring methodologies.

\section{6. Limitations}\label{limitations}

\begin{itemize}
\tightlist
\item
  \textbf{Data completeness:} Historical records are inherently
  incomplete. Communities that left no archival trace are absent from
  this dataset, introducing survivorship bias.
\item
  \textbf{Retrospective scoring:} Applying modern ESG categories to
  medieval institutions carries methodological risk. Scores reflect
  documented outputs, not intent.
\item
  \textbf{Geographic scope:} This study focuses on Western and Central
  Europe. Eastern European, Scandinavian, and Middle Eastern monastic
  traditions require separate treatment.
\item
  \textbf{Economic proxies:} Revenue indices are constructed from tithe
  records and land grants, which serve as proxies for actual economic
  output. Direct financial records are rare for this period.
\end{itemize}

\section{7. Conclusion}\label{conclusion}

Medieval monastic communities functioned as \textbf{integrated
institutional systems} managing environmental, economic, and governance
dimensions simultaneously.

The data supports three primary conclusions:

\begin{enumerate}
\def\labelenumi{\arabic{enumi}.}
\tightlist
\item
  \textbf{Governance protocol} (the Order's rule) is the strongest
  predictor of institutional longevity and ESG performance
\item
  \textbf{Territorial and economic diversification} correlate with
  resilience across geopolitical disruptions
\item
  \textbf{Heritage impact} --- architecture, agriculture, landscape ---
  represents an ultra-long-term outcome category that modern ESG
  frameworks do not yet capture
\end{enumerate}

These findings position medieval monastic history not as an antiquarian
subject, but as a \textbf{longitudinal laboratory for institutional
impact analysis}.

\section{References}\label{references}




\end{document}
